\section{Power energies, integral mixed inequalities, and capacity stability}
\label{sec:r11-power-energy-stability}

Keep the original arbitrary $S$, ordinary domain $V$, and $q=n-p\geq1$.
All top products are those of
Theorem~\ref{thm:r11-boundary-mixed-top-calculus}. The statements
below use integrated positivity, comparison, and fixed-domain
capacity. They do not assume a pointwise Alexandrov--Fenchel or
hyperbolicity property for the middle-positive background. Such a
property is not available for arbitrary $S$.

For a negative extended classical psh $u$ with $\mathrm{BG}(u,0)$
and $s>0$, define the extended nonnegative quantities
\[
 \begin{split}
 e_s(u)&=\int_V(-u)^s\,d\mu_u,\\
 J_s(u)&=\int_0^\infty t^{q+s-1}
                       \operatorname{Cap}_S(\{u<-t\};V)\,dt.
 \end{split}
 \tag{PE1}
\]
The first integral is over the finite-value locus; $\mu_u$ is zero
on the omitted pole set. The integrand defining $J_s$ is Borel
because its capacity factor is nonincreasing in $t$.

\begin{theorem}[A capacity-only criterion for finite power energy]
\label{thm:r11-power-energy-capacity-equivalence}
For $u$ as above and $s>0$, the following are equivalent:
\[
 J_s(u)<\infty
 \quad\Longleftrightarrow\quad
 u\in\mathcal F_{S,\partial}(V)\ \hbox{and}\ e_s(u)<\infty.
 \tag{PE2}
\]
On this class, more generally with extended values allowed after
membership in $\mathcal F_{S,\partial}$ is known,
\[
 {s\over2^{q+s}}J_s(u)
       \leq e_s(u)
       \leq {s2^q\over1-2^{-s}}J_s(u).
 \tag{PE3}
\]
Denote the equivalent class in (PE2) by
$\mathcal E^s_{S,\partial}(V)$. Raw or outer capacity gives the
same $J_s$, because all its sublevels are ordinary open.
\end{theorem}

\noindent\textbf{Status: \texttt{NEW\_PROVED}} for the
all-original arbitrary-background top calculus. Finiteness of
$e_s$ alone, outside the no-loss class, is not a sufficient
condition; the logarithmic counterexample has $\mu_u=0$.

\begin{proof}
For $u\in\mathcal F_{S,\partial}$, the first inequality of (BC4)
and Tonelli's theorem give
\[
 \begin{split}
 J_s(u)
 &\leq2^q\int_0^\infty t^{s-1}\mu_u(\{u<-t/2\})\,dt\\
 &= {2^{q+s}\over s}\int_V(-u)^s\,d\mu_u.
 \end{split}
 \tag{PE4}
\]
For the opposite direction no absence of loss is needed.
Real-level truncations are bounded classical psh functions and
obey the same consistency as integer truncations: one may express
them as the maximum of an integer truncation and a constant.
On $\{-2t<u<-t\}$ the measure $\mu_u$ equals the bounded
measure of $\max(u,-2t)$. Normalizing this potential to $[0,1]$
therefore gives
\[
       \mu_u(\{-2t<u<-t\})
          \leq(2t)^q\operatorname{Cap}_S(\{u<-t\};V).
 \tag{PE5}
\]
For a fixed finite value $a=-u>0$, the parameter $t$ in this
annulus ranges over $(a/2,a)$. Tonelli consequently yields
\[
 {1-2^{-s}\over s}e_s(u)
   =\int_0^\infty t^{s-1}
                 \mu_u(\{-2t<u<-t\})\,dt
   \leq2^qJ_s(u).
 \tag{PE6}
\]
Endpoint levels cause no loss in this integration, since for each
fixed value they form a Lebesgue-null set of parameters.

It remains to show that finite $J_s$ forces no loss. Write
$\kappa=q+s$. Monotonicity of the sublevel capacity implies
\[
 {1-2^{-\kappa}\over\kappa}\,
       t^\kappa\operatorname{Cap}_S(\{u<-t\};V)
 \leq\int_{t/2}^t r^{\kappa-1}
                 \operatorname{Cap}_S(\{u<-r\};V)\,dr.
 \tag{PE7}
\]
When $J_s<\infty$, this proves
$t^q\operatorname{Cap}_S(\{u<-t\};V)\to0$.
Theorem~\ref{thm:r11-capacity-tail-criterion} gives
$u\in\mathcal F_{S,\partial}$, and (PE6) gives finite $e_s$.
Conversely, (PE4) proves finiteness of $J_s$ from the right side
of (PE2). This proves every assertion without an assumption on
an unrestricted unbounded Monge--Ampere operator.
\end{proof}

Put
\[
 \kappa=q+s,\qquad
 A_{q,s}={2^{2q+s}\over1-2^{-s}},\qquad
 C_{q,s}={A_{q,s}(q+1)^{\kappa+1}\over q!}.
 \tag{PE8}
\]
These constants are explicit and are not claimed optimal.

\begin{theorem}[Power-energy cone and an integral mixed Holder inequality]
\label{thm:r11-power-mixed-holder}
For every $s>0$, the class $\mathcal E^s_{S,\partial}(V)$ is
upward closed among negative extended classical psh functions,
closed under finite positive sums, and closed under finite maxima.
If $u\leq v\leq0$ with $u$ in the class, then
\[
                         e_s(v)\leq2^{q+s}e_s(u).
 \tag{PE9}
\]
For $u_1,\ldots,u_m$ in the class,
\[
                  e_s\left(\sum_{i=1}^m u_i\right)
                    \leq A_{q,s}m^\kappa\sum_{i=1}^m e_s(u_i).
 \tag{PE10}
\]
For any $u_0,u_1,\ldots,u_q$ in the class, the full mixed
weighted integral is finite and satisfies
\[
 \boxed{\quad
 \int_V(-u_0)^s\,dM(u_1,\ldots,u_q)
 \leq C_{q,s}\,e_s(u_0)^{s/\kappa}
                         \prod_{i=1}^q e_s(u_i)^{1/\kappa}.
 \quad}
 \tag{PE11}
\]
A zero factor makes the right side zero. The mixed measure gives
zero mass to the pole of $u_0$, so the integral retains the
actual representative without any undefined product at a pole.
\end{theorem}

\noindent\textbf{Status: \texttt{NEW\_PROVED}}. This is an
integral Cegrell-type estimate with a nonsharp constant. It is
not a pointwise determinant inequality for $S$.

\begin{proof}
Upward membership follows from the solid boundary cone and
$J_s(v)\leq J_s(u)$. Integrating (BC12) by layer cake gives
the sharper numerical estimate (PE9). For a sum, (BC6) with
all coefficients one yields
\[
 J_s\left(\sum_i u_i\right)\leq m^\kappa\sum_iJ_s(u_i).
\]
Combine both inequalities of (PE3) to obtain (PE10).
Positive homogeneity $e_s(cu)=c^\kappa e_s(u)$ follows from
the exact scaling of the top measure. These facts prove the
stated cone properties; a maximum follows from upward closure.

Let $w=u_0+\cdots+u_q$. Its positive mixed expansion contains
$q!M(u_1,\ldots,u_q)$, hence
\[
                 q!M(u_1,\ldots,u_q)\leq\mu_w.
 \tag{PE12}
\]
Since $w=-\infty$ at every pole of $u_0$, the right side gives
zero mass to that pole, and so does the left side. On the
remaining locus, $(-u_0)^s\leq(-w)^s$. Thus the left side
of (PE11) is at most $e_s(w)/q!$ and is finite.

For $\varepsilon>0$ set
$v_i=u_i/(e_s(u_i)+\varepsilon)^{1/\kappa}$.
Then $e_s(v_i)\leq1$. By (PE10) with $m=q+1$,
$e_s(\sum_{i=0}^qv_i)/q!\leq C_{q,s}$.
Apply (PE12) to this normalized tuple and then use positive
homogeneity in all $q$ differentiated slots and the degree-$s$
homogeneity of the weight. This gives
\[
 \int(-u_0)^s\,dM(u_1,\ldots,u_q)
 \leq C_{q,s}(e_s(u_0)+\varepsilon)^{s/\kappa}
                       \prod_{i=1}^q(e_s(u_i)+\varepsilon)^{1/\kappa}.
\]
Letting $\varepsilon\downarrow0$ proves (PE11), including every
degenerate zero-energy case. No nondegeneracy of $S$ is used.
\end{proof}

\begin{proposition}[Capacity domination and quantitative top truncation tails]
\label{prop:r11-power-capacity-domination}
Let $u_1,\ldots,u_q\in\mathcal E^s_{S,\partial}(V)$, put
$w=\sum_i u_i$, and write $M=M(u_1,\ldots,u_q)$.
For every Borel $E\subset V$ and every $t>0$,
\[
 M(E)\leq{1\over q!}
       \left(q^qt^q\operatorname{Cap}_S(E;V)+t^{-s}e_s(w)\right).
 \tag{PE13}
\]
In particular, for finite-capacity sets,
\[
 M(E)\leq B_{q,s}
       \left(\prod_{i=1}^q e_s(u_i)^{1/\kappa}\right)
                    \operatorname{Cap}_S(E;V)^{s/\kappa},
 \quad
 B_{q,s}={q^q+1\over q!}
                  \left(A_{q,s}q^{\kappa+1}\right)^{q/\kappa}.
 \tag{PE14}
\]
For infinite capacity the estimate is vacuous. A zero energy
factor makes $M$ the zero measure.
If $M_k=T_{S,q}(\max(u_1,-k),\ldots,\max(u_q,-k))$, then
\[
 \|M_k-M\|_{\mathrm{TV}(V)}
       \leq {q^q4^{q+s}+1\over q!}\,k^{-s}e_s(w).
 \tag{PE15}
\]
All inequalities hold on actual Borel sets, not only on a chosen
moment image or a restricted family of test functions.
\end{proposition}

\noindent\textbf{Status: \texttt{NEW\_PROVED}}.

\begin{proof}
On $E_t=\bigcap_i\{u_i>-t\}$, strict locality identifies $M$
with the bounded mixed product of the real-level truncations.
Their sum lies in $[-qt,0]$. The positive expansion and capacity
normalization give
\[
                         M(E\cap E_t)
                        \leq{(qt)^q\over q!}\operatorname{Cap}_S(E;V).
\]
The complement of $E_t$ is contained in $\{w\leq-t\}$.
The positive expansion $q!M\leq\mu_w$ and the power moment
bound give $M(V\setminus E_t)\leq e_s(w)/(q!t^s)$.
This proves (PE13).

If $c=\operatorname{Cap}_S(E;V)$ and $e=e_s(w)$ are positive
and finite, choose $t=(e/c)^{1/\kappa}$ in (PE13), obtaining
\[
                   M(E)\leq{q^q+1\over q!}
                                   e^{q/\kappa}c^{s/\kappa}.
\]
The cases $c=0$ or $e=0$ follow by letting $t\to\infty$ or
$t\downarrow0$, respectively. Normalize the $q$ slots by
$(e_s(u_i)+\varepsilon)^{1/\kappa}$ as in the preceding proof;
(PE10) bounds the energy of their sum by
$A_{q,s}q^{\kappa+1}$. Scaling back and letting
$\varepsilon\downarrow0$ gives (PE14). If an energy is zero,
this proves vanishing on every compact set, whose relative
capacity is finite by frozen CLN; the Radon measure is therefore
zero.

For (PE15), use the exact agreement on $E_k$. On its complement,
the bounded positive expansion and (BC4) give
\[
 \begin{split}
 M_k(V\setminus E_k)
  &\leq{(qk)^q\over q!}\operatorname{Cap}_S(\{w<-k/2\};V)\\
  &\leq{q^q4^q\over q!}\mu_w(\{w<-k/4\})
   \leq{q^q4^{q+s}\over q!}k^{-s}e_s(w).
 \end{split}
\]
The limiting mass on the complement is at most
$k^{-s}e_s(w)/q!$. Their sum bounds total variation and proves
(PE15).
\end{proof}

\subsection{The exact capacity quantifier for nonmonotone convergence}

\begin{definition}[Truncated local convergence in original capacity]
\label{def:r11-truncated-local-capacity}
For negative extended classical psh functions $v_l,u$ on $V$,
write $v_l\to u$ in truncated local original capacity if the
following holds. For every ordinary analytic Laurent-box domain
$\Omega\Subset V$ whose compact closure is contained in an
original unit-moment chart, every compact $L\Subset\Omega$,
every $k>0$, and every $\delta>0$,
\[
 \operatorname{Cap}_{S|\Omega}
 \left(L\cap\{|\max(v_l,-k)-\max(u,-k)|>\delta\};\Omega\right)
                                               \longrightarrow0.
 \tag{PE16}
\]
The quantifier includes every actual finite refinement of that
moment chart and every further Laurent-box shrink. Under the
additional one-affinoid condition (DG1), Corollary
\ref{cor:r11-affinoid-local-capacity-topology} proves its equivalence
with fixed-domain capacity convergence on every compact core of $V$.
Without that domain hypothesis the printed all-local definition
remains the hypothesis used here.
\end{definition}

\begin{lemma}[The bounded finite-stage stability input on analytic boxes]
\label{lem:r11-bounded-analytic-capacity-stability}
For uniformly bounded classical psh families, (PE16) implies the
frozen weighted BT stability on every original compact test, in
every degree $0\leq r\leq q$. The same multiplier conclusion
holds for a bounded continuous scalar composition of one of
the potentials.
\end{lemma}

\noindent\textbf{Status: \texttt{DIRECT\_COROLLARY}} of the
frozen R1/R4 theorem \texttt{thm:capacity-stability}, with the
analytic bounded quasicontinuity of
Corollary~\ref{cor:r11-analytic-bounded-relative-QC} supplying
its domain extension. The bounded calculus itself is not new.

\begin{proof}
We record the precise controls, so no stronger topology is
silently substituted. Localize an original test to finitely many
of the stated boxes. Its finitely many moment coordinates $t_a$
are bounded on each compact closure. The bounded smooth classical
psh function $\rho=\frac12\sum_a t_a^2$ has
$D\rho=\sum_a d't_a\wedge d''t_a$. Finite original positive
matrix domination bounds every coefficient measure of a degree-$r$
bounded product on the compact stage by a constant times
\[
          T_{S,q}(a_1,\ldots,a_r,\rho,\ldots,\rho).
\]
Mixed capacity domination therefore bounds it by
$C\operatorname{Cap}_{S|\Omega}(\,\cdot\,;\Omega)$, uniformly
in the family. Only wedge with positive $(1,1)$ forms is used,
so the arbitrary middle-positive cone is retained.

The analytic Laurent collar theorem gives the exact bounded
quasicontinuity needed for the frozen capacity-controlled
multiplier lemma. Its proof separates a changing bounded weight
on its capacity-small gap set and replaces the fixed limit weight
by a continuous function off a small ordinary open exception.
The coefficient domination just given controls both errors.
Induction through $D(uP)$ then proves the bounded stability in
every degree, exactly as in the frozen proof. A bounded continuous
composition has capacity convergence by uniform continuity on
the bounded range and is quasicontinuous by composition, so the
same multiplier lemma applies to it. All controls in this proof use
the local capacity in (PE16). Under (DG1), the separate analytic
extension theorem (AL8)--(AL9) gives the equivalent fixed-domain
criterion; that comparison is not assumed in the multiplier proof.
\end{proof}

\begin{theorem}[Nonmonotone stability and lower semicontinuity of power energy]
\label{thm:r11-power-capacity-stability}
Fix $s>0$. For $1\leq i\leq q$, let
$v_{i,l}\in\mathcal E^s_{S,\partial}(V)$ and let $u_i$ be
negative extended classical psh with $\mathrm{BG}(u_i,0)$.
Assume (PE16) in every slot and
\[
                         \sup_{i,l}e_s(v_{i,l})<\infty.
 \tag{PE17}
\]
Then $u_i\in\mathcal E^s_{S,\partial}(V)$,
\[
 e_s(u_i)\leq\liminf_l e_s(v_{i,l}),\qquad
 M(v_{1,l},\ldots,v_{q,l})\longrightarrow M(u_1,\ldots,u_q)
 \tag{PE18}
\]
on every original compact smooth scalar test. No monotonicity of
the approximating potentials is assumed.

Suppose additionally that $v_{0,l}$ and $u_0$ satisfy the same
hypotheses. If $\Psi:[0,\infty)\to\mathbf R$ is continuous and
\[
                         |\Psi(t)|/t^s\longrightarrow0
                                      \quad(t\to\infty),
 \tag{PE19}
\]
then the locally finite weighted measures satisfy
\[
 \Psi(-v_{0,l})M(v_{1,l},\ldots,v_{q,l})
        \longrightarrow\Psi(-u_0)M(u_1,\ldots,u_q)
 \tag{PE20}
\]
on the same full original test family. Their weights at poles
are immaterial because the measures give those poles zero mass.
In particular powers $t^a$ with $0<a<s$ are allowed.
\end{theorem}

\noindent\textbf{Status: \texttt{NEW\_PROVED}}. The exact
localized capacity hypothesis and the global power-energy bound
are both retained.

\begin{proof}
First prove membership of a limit $u$, before applying any
unbounded continuity assertion. Fix $t,\delta>0$. The set
$\{u<-t-\delta\}$ has compact closure by its boundary-gap
condition. Choose a finite cover of that closure by compact
cores $L_a\Subset\Omega_a$ of admissible boxes. Choose
$k>t+\delta$. On the strict sublevel one has
\[
 \{u<-t-\delta\}
 \subset\{v_l<-t\}\cup
       \{|\max(v_l,-k)-\max(u,-k)|>\delta\}.
\]
The second set need only be considered on the fixed compact
cover. Domain restriction and (PE16) make its contribution to
the capacity measured on $V$ tend to zero. Consequently
\[
 \operatorname{Cap}_S(\{u<-t-\delta\};V)
       \leq\liminf_l\operatorname{Cap}_S(\{v_l<-t\};V).
\]
Let $\delta\downarrow0$. The raw capacity is continuous from
below: this follows by commuting the two increasing suprema
over capacity testers and over the increasing carriers. Hence
the same inequality holds with $\{u<-t\}$ on the left.
Fatou and (PE4) give
\[
 J_s(u)\leq\liminf_lJ_s(v_l)
       \leq {2^{q+s}\over s}\liminf_l e_s(v_l)<\infty.
 \tag{PE21}
\]
Theorem~\ref{thm:r11-power-energy-capacity-equivalence} proves
that every $u_i$ lies in the energy class. This step used only
bounded capacity and the printed boundary condition.

Put $w_l=\sum_i v_{i,l}$ and $w=\sum_i u_i$.
Equation (PE10) gives one finite uniform bound for $e_s(w_l)$
and $e_s(w)$. By (PE15), replacing every differentiated slot
by its truncation at $-k$ changes its mixed top measure by at
most a constant times $k^{-s}$ in total variation, uniformly
in $l$. For fixed $k$, Lemma
\ref{lem:r11-bounded-analytic-capacity-stability} gives convergence
of the bounded products on every original compact test.
Pass first in $l$ and then in $k$; the explicit uniform
variation bound proves the unweighted convergence in (PE18).

We next prove (PE20). The integral inequality (PE11) and the
uniform energy bound give
\[
 \sup_l\int(-v_{0,l})^s\,dM(v_{1,l},\ldots,v_{q,l})<\infty,
 \qquad
 \int(-u_0)^s\,dM(u_1,\ldots,u_q)<\infty.
 \tag{PE22}
\]
For $R>0$ let $\Psi_R(t)=\Psi(\min(t,R))$ and put
$\epsilon_R=\sup_{t\geq R}|\Psi(t)|/t^s$ for large $R$.
Then $\epsilon_R\to0$ and
\[
          |\Psi(t)-\Psi_R(t)|
                \leq2\epsilon_R t^s\mathbf1_{\{t>R\}}.
\]
By (PE22), the errors from replacing the weight by $\Psi_R$
are uniformly small on every compact test as $R\to\infty$.
For fixed $R$, the weights are bounded continuous compositions
of truncations of the zeroth potential. Replacing the other
slots by their truncations at $-k$ has a uniform error bounded
by $\|\Psi_R\|_\infty$ times (PE15). For fixed $R,k$, the
bounded multiplier statement of the preceding lemma passes
the limit in $l$. Let $k\to\infty$ and finally $R\to\infty$.
This proves (PE20). Local finiteness follows from (PE22) on
the large-value region and the uniform finite compact masses
on the remaining bounded-value region. The positive expansion
with the zeroth potential proves that its pole has zero mixed
mass, as in (PE12).

Finally apply the bounded-weight case just proved to a single
potential repeated in every differentiated slot, with weight
$\min(t^s,R^s)$ and any nonnegative original compact smooth
$0\leq\chi\leq1$. It gives
\[
 \int\chi\min((-u)^s,R^s)\,d\mu_u
  =\lim_l\int\chi\min((-v_l)^s,R^s)\,d\mu_{v_l}
                         \leq\liminf_l e_s(v_l).
\]
Increase $R$ and then take the supremum over such compact
cutoffs. Radon inner regularity gives $e_s(u)$ on the left.
This proves the sharp lower-semicontinuity assertion in (PE18),
improving the preliminary bound from (PE21).
\end{proof}

\begin{corollary}[A precise closure principle for class-valued nonlinear Dirichlet data]
\label{cor:r11-power-Dirichlet-closure}
Fix $s>0$ and a positive Radon datum $\mu$ on $V$. Suppose there
are bounded negative classical psh $u_l$ satisfying
$\mathrm{BG}(u_l,0)$,
\[
 (Du_l)^q\wedge S=\mu_l,\qquad
 \mu_l\longrightarrow\mu\ \hbox{on all original compact tests},
 \qquad\sup_l e_s(u_l)<\infty,
 \tag{PE23}
\]
and suppose $u_l$ has a limit $u$ in the exact sense (PE16),
where $u$ is actual negative extended classical psh and
$\mathrm{BG}(u,0)$. Then
\[
                    u\in\mathcal E^s_{S,\partial}(V),
                    \qquad\mu_u=\mu.
 \tag{PE24}
\]
Any two solutions in this class with datum $\mu$ agree outside
a relative outer-capacity-null set.
\end{corollary}

\noindent\textbf{Status: \texttt{CONDITIONAL}} as a nonlinear
Dirichlet existence mechanism: bounded solvability, the uniform
power-energy estimate, an actual classical capacity limit, and
its zero-boundary gap condition are the explicit premises.
The passage to the limit and the stated uniqueness conclusion
are fully proved. No existence or compactness premise is declared
automatic for arbitrary original $S$.

\begin{proof}
Apply Theorem~\ref{thm:r11-power-capacity-stability} with the
same potential in every slot. The original current limit is
both $\mu_u$ and $\mu$, proving (PE24). Outer-capacity uniqueness
is Corollary~\ref{cor:r11-boundary-outer-domination}.
\end{proof}

\paragraph{Precise remaining gates.}
The theory above supplies full-dimensional arbitrary-background
energy inequalities and stability in a specified unbounded
boundary class. It does not prove capacity compactness of an
arbitrary energy-bounded family, regularization of a proposed
unbounded envelope, pointwise uniqueness for degenerate $S$,
or unrestricted nonlinear Dirichlet existence. The original
all-refinement PS/PL quotient and CTO remain separate open gates.
